\documentclass[11pt,a4paper]{article}
\usepackage[spanish,es-nodecimaldot]{babel}
\usepackage{fontspec}
\setsansfont{DejaVu Sans}
\renewcommand{\familydefault}{\sfdefault}
\usepackage[margin=1.8cm]{geometry}
\usepackage{enumitem,xcolor,tabularx,array,booktabs,needspace}
\usepackage{hyperref}
\definecolor{gris}{HTML}{555B63}\definecolor{azul}{HTML}{224765}\definecolor{fondo}{HTML}{F5F7FA}
\hypersetup{colorlinks=true,urlcolor=azul,linkcolor=azul}
\setlist[enumerate]{leftmargin=*,itemsep=4pt,topsep=3pt}
\setlist[itemize]{leftmargin=*,itemsep=2pt,topsep=2pt}
\setlength{\parindent}{0pt}
\setlength{\parskip}{3pt}
% Campo de texto rellenable; al imprimir queda como recuadro vacío.
% #1 nombre único, #2 alto
\newcommand{\campo}[2]{\par\nopagebreak\noindent\TextField[name=#1,multiline=true,width=\linewidth,height=#2,bordercolor={0.6 0.6 0.6},backgroundcolor={0.96 0.97 0.98},charsize=9pt]{}\par\vspace{4pt}}
\newcommand{\linea}[2]{\TextField[name=#1,width=#2,height=13pt,bordercolor={0.6 0.6 0.6},backgroundcolor={0.96 0.97 0.98},charsize=9pt]{}}
\newcommand{\nivel}[1]{\needspace{8\baselineskip}\section*{\color{azul}#1}}
\newcommand{\caja}[1]{\par\vspace{2pt}\noindent\fcolorbox{gris}{white}{\parbox{0.95\linewidth}{\small #1}}\par\vspace{4pt}}
\begin{document}
\begin{Form}
\begin{center}
{\LARGE\bfseries Sesión 2: misiones con IA}\\[5pt]
{\small Ficha digital · IA Total · San Juan Bosco}
\end{center}
\textbf{Nombre:} \linea{nombre}{7.5cm}\hfill\textbf{Fecha:} \linea{fecha}{3.5cm}\\[4pt]
\textbf{Proyecto final:} \linea{proyecto}{8.4cm}\hfill\textbf{Equipo:} \linea{equipo}{3.5cm}

\smallskip
\textcolor{gris}{\small \textbf{Cómo usarla:} ábrela en el navegador (Chrome, Edge o Firefox) o en Acrobat, escribe en los recuadros y guárdala con «Descargar» o «Guardar». También puede imprimirse. Copia aquí tus prompts y observaciones breves, no respuestas largas. No escribas datos personales en la IA.}

\nivel{Calentamiento · ¿Qué palabra sigue?}
Escribe la primera palabra que se te ocurra para completar cada frase. Tienes 10 segundos por ronda.

1. «En la selva amazónica viven muchos…» \linea{jr1}{5cm}\\[3pt]
2. «Mi comida favorita es el…» \linea{jr2}{5cm}\\[3pt]
3. «La primera panadería de San Juan Bosco la abrió…» \linea{jr3}{5cm}\\[3pt]
¿Sabías de verdad la respuesta de la ronda 3? ¿Qué tiene que ver con cómo escribe una IA?
\campo{jidea}{1.2cm}

\nivel{Fórmula y plantilla de un buen prompt}
\textbf{Rol + Tarea + Contexto + Formato + Tono + Límite}

\small
\noindent\begin{tabularx}{\linewidth}{@{}p{2.4cm}X@{}}\toprule
\textbf{Rol} & ¿Desde qué perspectiva debe responder?\\
\textbf{Tarea} & ¿Qué debe hacer exactamente?\\
\textbf{Contexto} & ¿Qué información necesita saber?\\
\textbf{Formato} & ¿Cómo quiero recibir la respuesta? (correo, tabla, lista, informe)\\
\textbf{Tono} & ¿Formal, amable, directo, institucional?\\
\textbf{Límite} & ¿Cuántas palabras, puntos o párrafos? ¿Qué no debe inventar?\\\bottomrule
\end{tabularx}
\normalsize

\textbf{Plantilla para copiar} (úsala en todos los niveles):
\caja{Actúa como [rol]. Necesito que [tarea]. El contexto es [información clave]. El resultado debe estar en formato [correo / tabla / lista / informe]. Usa un tono [formal / amable / institucional / directo]. Limita la respuesta a [número de palabras / párrafos / puntos].}
\textbf{Pruébala ahora.} Copia este prompt junto con las notas y pégalo en la IA:
\caja{«Actúa como coordinador de un club escolar. Necesito que conviertas estas notas en una lista de tareas. El contexto es una reunión para organizar nuestro puesto en la Feria Comunitaria. El resultado debe estar en una tabla. Usa un tono claro y amable. Limita cada tarea a una frase.»\\[3pt]
\textbf{Notas:} conseguir dos mesas; preguntar al comité por extensiones eléctricas; preparar carteles del proyecto; llegar una hora antes; decidir quién explica el proyecto a los visitantes.}
¿La IA respetó los seis elementos? ¿Cuál cambiarías?
\campo{plantobs}{1.2cm}

\nivel{Nivel 1 · Misma pregunta, distintas IA}
Envía esta petición a \textbf{al menos dos} asistentes (ChatGPT, Gemini o Claude):
\caja{«Explica qué es la biodiversidad a un estudiante de 15 años en tres frases y da un ejemplo de la Amazonía.»}
¿Qué asistentes usaste? \linea{n1ia}{10.5cm}\\[3pt]
¿Cuál respuesta fue más clara y por qué?
\campo{n1clara}{1.6cm}
¿Hay algún dato que deberías comprobar? ¿Cuál?
\campo{n1dato}{1.2cm}

\nivel{Nivel 2 · Dale un rol}
Escribe tu propio prompt para \textbf{dos} de estas tareas (la tercera es reto extra). Indica el rol al inicio.
\begin{itemize}
\item[\textbf{A.}] \textbf{Guía turístico:} un plan de un día para jóvenes que visitan San Juan Bosco. Pide que escriba [lugar] cuando no conozca un sitio concreto, en vez de inventarlo.
\item[\textbf{B.}] \textbf{Profesor de colegio:} explicar la erosión del suelo con un ejemplo de agricultura.
\item[\textbf{C.}] \textbf{Contador:} explicar a una familia que vende queso cómo calcular su ganancia mensual, con números ficticios.
\end{itemize}
Prompt 1 (letra y texto):
\campo{n2p1}{1.6cm}
Prompt 2 (letra y texto):
\campo{n2p2}{1.6cm}
¿Qué cambió en la respuesta al darle un rol?
\campo{n2cambio}{1.2cm}

\nivel{Nivel 3 · Formato, tono y límite}
Escribe tu propio prompt para \textbf{dos} de estas tareas (la tercera es reto extra). Pide un formato, un tono y \textbf{al menos un límite}: número de palabras o puntos, o «no inventes datos».
\begin{itemize}
\item[\textbf{A.}] \textbf{Tabla:} tres actividades de un pequeño negocio familiar, una dificultad de cada una y una posible mejora.
\item[\textbf{B.}] \textbf{Mini plan de negocio:} vender jugos naturales en la Feria, en cinco puntos: producto, clientes, costos, precio y un riesgo.
\item[\textbf{C.}] \textbf{Resumen en cuatro ideas} del texto siguiente, sin añadir información.
\end{itemize}
\caja{Durante una actividad escolar, un grupo de estudiantes observó residuos cerca de un río. Decidieron separar los materiales que podían reutilizarse y conversar con sus familias sobre cómo reducir la basura. También propusieron colocar recipientes identificados en la escuela y organizar una jornada de limpieza con personas voluntarias. Antes de realizarla, acordaron consultar a docentes y autoridades sobre los permisos y la disposición segura de los residuos. El grupo concluyó que cuidar el agua requiere acciones constantes, no solo una jornada.}
Prompt 1 (letra y texto):
\campo{n3p1}{1.6cm}
Prompt 2 (letra y texto):
\campo{n3p2}{1.6cm}
¿La IA respetó el formato, el tono y el límite? ¿Añadió algo que no pediste?
\campo{n3revision}{1.2cm}

\nivel{Nivel 4 · Conversa con la IA}
La primera respuesta es un borrador. En el mismo chat de una respuesta de los niveles 2 o 3, envía estos tres mensajes, uno por uno:
\caja{1. «Hazlo más corto y sencillo.»\quad 2. «Agrega un ejemplo de San Juan Bosco.»\quad 3. «¿Qué datos de tu respuesta debo comprobar?»}
¿Qué respuesta elegiste? \linea{n4cual}{10cm}\\[3pt]
¿Mejoró? ¿El ejemplo de San Juan Bosco es real o hay que comprobarlo?
\campo{n4mejora}{1.4cm}
\textbf{Reto extra · Dale un modelo.} Pega un texto corto que te guste y pide: «Escribe otro texto sobre [tema] con el mismo estilo que este: …». Tu prompt:
\campo{n4modelo}{1.2cm}

\nivel{Batalla de prompts}
\textbf{Objetivo:} un mensaje de \textbf{máximo 50 palabras} que invite a jóvenes de otra ciudad a visitar San Juan Bosco, \textbf{sin inventar lugares ni datos}. Tienes 6 minutos. Gana el mensaje más claro y atractivo que cumpla las reglas.

Tu mejor prompt:
\campo{batprompt}{1.8cm}
El mensaje que obtuviste:
\campo{batmensaje}{1.8cm}

\nivel{Cazadores de alucinaciones}
Pide a la IA tres datos concretos de San Juan Bosco o Morona Santiago (año de cantonización, habitantes, ríos, parroquias…) y compruébalos en una fuente confiable: sitio oficial del municipio, INEC u otra institución.

\small
\noindent\begin{tabularx}{\linewidth}{@{}p{3.4cm}XXX@{}}
\toprule
\textbf{Dato pedido} & \textbf{Lo que dijo la IA} & \textbf{Lo que dice la fuente} & \textbf{Fuente (enlace)}\\\midrule
\linea{al1a}{3.1cm} & \linea{al1b}{\linewidth} & \linea{al1c}{\linewidth} & \linea{al1d}{\linewidth}\\[4pt]
\linea{al2a}{3.1cm} & \linea{al2b}{\linewidth} & \linea{al2c}{\linewidth} & \linea{al2d}{\linewidth}\\[4pt]
\linea{al3a}{3.1cm} & \linea{al3b}{\linewidth} & \linea{al3c}{\linewidth} & \linea{al3d}{\linewidth}\\\bottomrule
\end{tabularx}
\normalsize

\medskip
\needspace{4\baselineskip}¿Encontraste un dato equivocado o inventado? ¿Cuál?
\campo{alerror}{1.2cm}

\nivel{Estaciones de ofimática}
¿Cómo trabajaste? \linea{modalidad}{12.3cm}\\
\textcolor{gris}{\small Por ejemplo: «Gemini dentro de Gmail y Docs», «ChatGPT + Google Docs» o «ChatGPT + Word».}

\subsection*{1. Correo}
Redacta con IA un correo formal que solicite información a una institución (municipio, SENESCYT o una empresa local). \textbf{Guárdalo como borrador; no lo envíes.}

Institución y asunto del correo:
\campo{correoasunto}{1.0cm}
Después, pide a la IA un resumen en tres puntos de este hilo de correos ficticio:
\caja{\textbf{Mensaje 1 (Comité de la Feria):} Hola a todos: la Feria Comunitaria será en la cancha de la escuela. Necesitamos saber cuántas mesas usará cada proyecto antes del viernes.\\[3pt]
\textbf{Mensaje 2 (Estudiante):} Nuestro proyecto necesita una mesa y una conexión eléctrica para el computador. ¿Habrá extensiones?\\[3pt]
\textbf{Mensaje 3 (Comité de la Feria):} Habrá dos extensiones por fila; si necesitan más, tráiganlas. Lleguen una hora antes para armar los puestos. La fecha y el horario se confirmarán la próxima semana.}
Resumen obtenido (¿inventó alguna fecha?):
\campo{correoresumen}{1.6cm}

\subsection*{2. Documento}
En un documento, escribe con IA una invitación breve a la Feria Comunitaria que incluya esta idea: «Invitamos a las familias a conocer los proyectos preparados por estudiantes de nuestra escuela». Pide la traducción al inglés y comprueba que conserva las ideas de \emph{familias}, \emph{proyectos} y \emph{escuela}.

Invitación en español:
\campo{docesp}{1.6cm}
Traducción al inglés:
\campo{docing}{1.6cm}
Una corrección que hiciste:
\campo{doccorr}{1.0cm}

\subsection*{3. Hoja de cálculo}
Copia esta tabla en Google Sheets o Excel (cifras ficticias). Pide a la IA las fórmulas para el subtotal, el total de gastos y la ganancia; luego crea un gráfico de los gastos.

\begin{center}\small
\begin{tabular}{@{}lrrl@{}}\toprule
\textbf{Producto} & \textbf{Cantidad} & \textbf{Precio unitario} & \textbf{Subtotal}\\\midrule
Naranjas & 30 & \$0,15 & ?\\
Azúcar (kg) & 2 & \$1,10 & ?\\
Vasos & 50 & \$0,04 & ?\\
Hielo (fundas) & 2 & \$1,00 & ?\\\midrule
\textbf{Total de gastos} & & & ?\\
\textbf{Ingresos:} 40 jugos vendidos & 40 & \$0,50 & ?\\
\textbf{Ganancia} & & & ?\\\bottomrule
\end{tabular}
\end{center}
Fórmulas que usaste:
\campo{hojaformulas}{1.4cm}
Total de gastos: \linea{hojatotal}{3cm}\quad Ganancia: \linea{hojagan}{3cm}\\[3pt]
¿Comprobaste el cálculo en la hoja? \linea{hojacomp}{8cm}


\nivel{Avance del proyecto final (en equipo)}
Siéntense juntos; cada integrante trabaja en su computador con un rol. Al final, \textbf{todos copian en su ficha la versión acordada por el equipo}.

\small
\noindent\begin{tabularx}{\linewidth}{@{}p{3.2cm}X@{}}\toprule
\textbf{Coordinador} & Reúne las notas de todos, consolida la versión final y une las diapositivas del paso 3.\\
\textbf{Organizador} & Ejecuta el prompt del paso 1.\\
\textbf{Jurado} & Pide a la IA las preguntas difíciles del paso 2.\\
\textbf{Entrevistador} & Pide a la IA las preguntas para la comunidad del paso 2.\\
\textbf{Verificador} & Compara lo que dio la IA con las notas: ¿inventó algo? (En equipos de 4, lo hace el coordinador.)\\\bottomrule
\end{tabularx}
\normalsize

Mi rol: \linea{prol}{6cm}

\textbf{Paso 1 · Ordenar.} Copia este prompt, completa los corchetes y pégalo en la IA:
\caja{«Actúa como asesor de proyectos escolares. Nuestro proyecto es: [describe la idea]. Esta es la información que recolectamos: [pega tus notas]. Organízala en cuatro partes: problema, a quién afecta, qué sabemos y qué nos falta averiguar. Usa un tono claro y limita cada parte a dos frases. No inventes datos; si algo no está en mis notas, ponlo en "qué nos falta averiguar".»}
Revisa la respuesta, corrígela y cópiala aquí:

\small
\noindent\begin{tabularx}{\linewidth}{@{}p{3.6cm}X@{}}
\textbf{Problema} & \campo{pproblema}{1.2cm}\\
\textbf{A quién afecta} & \campo{pafecta}{1.2cm}\\
\textbf{Qué sabemos} (y de dónde lo sabemos) & \campo{psabemos}{1.8cm}\\
\textbf{Qué nos falta averiguar} & \campo{pfalta}{1.4cm}\\
\end{tabularx}
\normalsize

\textbf{Paso 2 · Poner a prueba.} El jurado pide a la IA: «Hazme cinco preguntas difíciles que un jurado haría sobre este proyecto»; el entrevistador: «Propón cinco preguntas para entrevistar a las personas afectadas». Luego el equipo elige y responde.

Responde aquí dos preguntas del jurado:
\campo{pjurado}{2.0cm}
Las tres mejores preguntas para entrevistar:
\campo{pentrevista}{1.6cm}

\textbf{Paso 3 · Presentar.} Preparen cuatro diapositivas (Google Slides, PowerPoint o Impress), \textbf{una por integrante}, con ayuda de la IA: \emph{1. Problema · 2. A quién afecta · 3. Lo que averiguamos · 4. Próximo paso}. Pocas palabras y solo datos que el equipo comprobó. El coordinador las une; en equipos de cuatro, también hace una. Las usarán en la presentación de 90 segundos.

Mi diapositiva (número y texto):
\campo{pdiapo}{1.4cm}

\textbf{Paso 4 · Intercambio.} Lo que les dijo el otro equipo:

\noindent\begin{tabularx}{\linewidth}{@{}XX@{}}
\textbf{Una estrella} (algo que funciona) & \textbf{Un deseo} (algo que mejorar)\\
\campo{pestrella}{1.2cm} & \campo{pdeseo}{1.2cm}\\
\end{tabularx}

\textbf{Próximo paso antes de la sesión 3:}
\campo{pproximo}{1.2cm}

\nivel{Misiones de reserva}
\textcolor{gris}{\small Solo si el instructor las activa.}

\textbf{Prompt inverso.} Escribe el prompt que crees que produjo la respuesta proyectada y pruébalo:
\campo{rinverso}{1.4cm}
\textbf{Que la IA les haga preguntas} (en equipo, con el prompt de la diapositiva). Una pregunta que no supieron responder:
\campo{rpregunta}{1.0cm}

\nivel{Boleto de salida}
Una cosa en la que la IA me ayuda:
\campo{bayuda}{0.9cm}
Una precaución que tendré:
\campo{bprecaucion}{0.9cm}
Una duda que todavía tengo:
\campo{bduda}{0.9cm}

\needspace{10\baselineskip}\subsection*{Autoevaluación}
\textcolor{gris}{\small Responde: sí, más o menos o no.}

\small
\noindent\begin{tabularx}{\linewidth}{@{}X l@{}}
Escribí prompts con rol, tarea, contexto, formato, tono y límite. & \linea{ae1}{3.2cm}\\[4pt]
Comprobé al menos un dato en una fuente confiable. & \linea{ae2}{3.2cm}\\[4pt]
Revisé y corregí lo que la IA escribió antes de usarlo. & \linea{ae3}{3.2cm}\\[4pt]
Mejoré una respuesta conversando con la IA. & \linea{ae5}{3.2cm}\\[4pt]
Cumplí mi rol en el avance del proyecto. & \linea{ae4}{3.2cm}\\
\end{tabularx}
\normalsize

\subsection*{Tarea para la próxima semana}
\begin{enumerate}
\item Usa una IA para preparar o mejorar una tarea escolar real. Anota qué le pediste, qué cambiaste y qué comprobaste.
\item Avanza en tu proyecto: consigue la información que quedó en «qué nos falta averiguar».
\end{enumerate}
Guarda esta ficha y entrégala por el medio que indique el instructor.
\end{Form}
\end{document}
